
\label{sec:instruction_following}

In the next set of experiments, we study how well \MethodName{} can follow language instructions, including its ability to carry out a variety of different contexts and follow referential instructions that differ systematically from the training data. Our experiments focus specifically on performing tasks in messy environments where there are many possible tasks that the robot could perform, requiring \MethodName{} to pay careful attention to the provided instruction in order to succeed. We find that \MethodName{} displays instruction following capabilities that significantly improve upon our prior models \Pizf{}~\cite{black2025pi05} and \Pizs~\cite{pi06model}.

\noindent \textbf{\MethodName{} can be flexibly prompted to perform a wide variety of different tasks.} Language following has presented a notorious challenge for robotic foundation models, particularly with open-vocabulary instructions that do not correspond directly to ones seen in training. In these experiments, we aim to study the breadth of \MethodName{}'s capabilities, to answer: can \MethodName{} better handle a wider variety of language instructions than prior models?

We evaluate \MethodName{} on a diverse set of instructions across 4 unseen kitchens and 2 unseen bedrooms that were not present in the training data (Fig.~\ref{fig:instruction_following}). Each experiment tests whether the robot can follow a 3 to 6 step sequence of instructions to achieve a specific goal. The tasks require a variety of different realistic tasks in kitchen and bedroom environments, including re-arranging and tidying items, interacting with furniture, and cleaning up spills. This combination of new test environments and diverse instructions presents a major challenge for robotic foundation models, which can struggle to follow even simple instructions in seen environments. We find that \MethodName{} is able to significantly outperform \Pizf{} and \Pizs{}, with a high overall instruction following success rate.

\noindent \textbf{\MethodName{} can handle out-of-distribution referential instructions. }
Because of the breadth and diversity of our training data, it is generally difficult for us to quantify how \emph{novel} the test instructions are. In the next set of experiments, we intentionally designed a set of instructions that are unusual, refer to objects in unconventional ways, or require understanding spatial relations. In Fig.~\ref{fig:instruction_generalization}, we compare \MethodName{} and prior models on a set of object re-arrangement instructions, broken up into \emph{standard} and \emph{complex} instructions. The standard instructions are phrased in a similar way to language instructions in the training data. The complex instructions use unusual language or complex spatial references, such as ``pick up an object I would use to eat soup" or ``pick up the fruit on the largest plate." In these experiments, we also evaluate \MethodName{} with and without subgoal images, which are generated by a lightweight world model as discussed previously. We can see that \MethodName{} improves over prior models on the complex instructions, and the use of subgoal images further boosts its performance, importing semantic understanding from the world model.

\begin{figure*}[h!]
    \includegraphics[width=\textwidth]{compressed_images_final/cross_embodiment.pdf}
    \caption{\textbf{Cross-embodiment transfer: }Left: Both \MethodName{} and prior models achieve strong cross-embodiment transfer directly out of the box on simpler re-arrangement or repositioning style tasks. For example, for the ``Table Setting'' task, data was collected with a variety of different robots, and the task was tested on the static bimanual robot. All of the models performed well. For tasks that required transfer from the bimanual UR5e robot to the smaller static bimanual robot (``Bag In Backpack'' and ``Organize Tupperware''), the embodiment gap is larger, as there is only one source robot that is bigger and heavier than the target robot. Here, \Pizf{} performed very poorly, but \Pizs{} still performs quite well. Transferring from the smaller static bimanual robot to the UR5e for the ``Shirt Bagging'' task introduces the largest embodiment gap, and here the \MethodName{} model significantly outperforms prior models. Right: for the more dexterous tasks that require folding towels and t-shirts, the embodiment gap poses an even greater challenge. Data for these tasks was collected with the smaller static bimanual robot, and the task was evaluated on the larger bimanual UR5e platform. \MethodName{} was able to transfer this task successfully, and improved even more when using visual subgoal images generated by our lightweight world model. In fact, task progress matches the 
    ``zero-shot'' performance of our most experienced human teleoperators, who have operated a diverse set of robots and attempted this task on the UR5e for the first time.}
    \label{fig:cross_embodiment}
\vspace{-5pt}
\end{figure*}

\noindent \textbf{\MethodName{} can follow instructions that go against dataset biases. }
Dataset bias presents a major challenge for instruction following: if the robot always does the same thing in a given scene, a model trained on such data will often ignore language in these scenes, blindly copying the behavior seen in the data. In the next experiment, we construct scenarios that suffer from this problem, and test if we can prompt \MethodName{} to go \emph{against} the natural bias in the dataset.
We constructed two tasks: ``Reverse Bussing'' and ``Reverse Fridge to Microwave.'' In our dataset, the ``bussing'' task involves putting trash in a trash bin and dishes in a bussing bin. For the ``Reverse Bussing'' task, the robot is asked to do the opposite: put the trash in the bussing bin, and the dishes in the trash. The ``Fridge to Microwave'' task requires taking food out of the fridge and putting it in the microwave, and we did not collect data going the other way. At test time, in the ``reverse'' version of this task (``Reverse Fridge to Microwave''), we prompt the robot to take the food from the microwave to the fridge, violating the bias in the dataset.

The results in Fig.~\ref{fig:compositional_generalization} show that \MethodName{} significantly improves over prior models on these tasks. This suggests that \MethodName{} has significantly better language following capabilities, and pays enough attention to the instructions to overcome the bias in the data for these tasks. On the ``Reverse Fridge to Microwave'' task, conditioning on generated subgoal images (\MethodName{} (GC)) is critical for success, as the world model can generate subgoals based on textual instructions effectively by leveraging web-scale image generation pre-training.


